\section{Reusable domain libraries}
\label{sec:domains}

The framework is most useful when it becomes a small collection of shared
object-language libraries.  This section describes the interfaces and the
proof reductions they support.

\subsection{Finite algebra and polynomial certificates}

A multivariate integer polynomial is a finite list of coefficient/exponent
pairs in a canonical order.  Addition and multiplication merge and convolve
those lists.  Normalization combines equal monomials and removes zero
coefficients.  Evaluation is a finite fold.

\begin{proposition}[Polynomial certificate reflection]
\label{prop:poly-reflection}
There is a primitive-recursive checker
\(\mathsf{PolyId}(P,Q,c)\) such that \(\ipra\) proves: if \(c\) is an accepted
normalization certificate for \(P-Q\), then
\[
 \forall \mathbf x\quad P(\mathbf x)=Q(\mathbf x).
\]
The same method checks rational identities after multiplying by a certified
positive common denominator.
\end{proposition}

\begin{proof}
The certificate lists every normalization step: sorting, combining equal
monomials, and deleting zeros.  The checker verifies local list equations.
Induction on the certificate length proves that evaluation is preserved.
Acceptance says the final list is empty, whose evaluation is zero.  The
induction invariant is equality of two integer evaluations and is
quantifier-free.  Variables \(\mathbf x\) remain free throughout and are
universally introduced at the end.
\end{proof}

This library handles ring identities, determinant expansions, adjugate
identities, trigonometric polynomial identities in \(\sin\) and \(\cos\), and
the finite coefficient parts of power-series proofs.  It also separates two
checks that are often conflated: a computer algebra system may produce the
certificate, while the small arithmetic checker proves its soundness.

\paragraph{Algebraic roots.}
A root should be represented first by an isolating rectangle or disk, a root
count, and a refinement procedure.  Polynomial arithmetic, interval bounds,
and a finite argument-principle or subdivision certificate validate each
region.  Repeated roots are represented by clusters with multiplicity data.
An ordered list of individual roots is added only under pairwise separation.
This resolves the branch discontinuity at collisions without weakening the
algebraic theorem.

\subsection{Linear algebra and geometry}

Matrices over coded rationals are finite lists.  Matrix multiplication,
transpose, determinant, minors, adjugate, dot products, and cross products are
finite polynomial operations.  Fixed-dimension identities are direct
polynomial certificates; variable-dimension identities use fold invariants.

\paragraph{Nonsingular systems.}
Cramer's rule and adjugate inversion are constructive under a determinant
apartness witness.  The transducer carries a lower bound on
\(\abs{\det A}\), requests sufficient entry precision, and proves the inverse
error by \cref{lem:reciprocal,lem:product-error}.  Gaussian elimination is
also possible, but every pivot choice must use a robust lower-bound scan.

\paragraph{Singular systems and eigenvectors.}
An exact nonzero kernel vector is not a continuous function of a singular
matrix near rank changes.  The unconditional output should be one of:
\begin{itemize}
\item a projective kernel certificate;
\item a residual-small unit vector at requested precision;
\item an exact normalized vector under a supplied minor or spectral gap;
\item a coherent isolating subspace projector rather than a basis.
\end{itemize}
Finite-net minimization of \(\norm{Av}^2\) on a rational sphere chart gives
the approximate version.  A gap promotes it to a stable exact name.

\paragraph{Euclidean and projective geometry.}
Points, lines, circles, conics, and incidences are homogeneous coordinate
codes.  Classical incidence theorems often reduce to determinant or polynomial
identities after nondegeneracy denominators are recorded.  Power-of-a-point,
Pappus, Ceva, Menelaus, and many circle theorems fit this pattern.  The library
should expose determinant certificates, signed-area predicates, and chart
transitions, not a large opaque geometry ontology.

\paragraph{Trigonometry and spherical geometry.}
Use regular names for \(\sin\) and \(\cos\) built from finite series or a
circle parameter, and state identities first as polynomial equations in a
pair \((c,s)\) satisfying \(c^2+s^2=1\).  Division-based tangent or half-angle
forms require denominator gaps.  Spherical theorems should prefer vector dot
and triple products; inverse angles are added via overlapping charts and
separation certificates.

\subsection{Real and complex analysis}

The shared analytic library should have four layers.

\begin{enumerate}
\item \textbf{Regular names and interval arithmetic.}  This is the numerical
      kernel of \cref{sec:represented}.
\item \textbf{Finite sums, products, and series tails.}  Every named special
      function exports a truncation algorithm, magnitude bound, and tail
      modulus.
\item \textbf{Mesh calculus.}  Integration, differentiation, and changes of
      variable use finite mesh sums plus local remainder moduli.
\item \textbf{Local analytic data.}  A holomorphic chart contains coefficient
      names and a majorant on a larger rational disk.  Cauchy estimates and
      termwise operations become finite tail calculations.
\end{enumerate}

A contour is a finite polygonal or piecewise rationally parametrized path.  A
residue theorem instance needs a positive distance from the path to every
pole, isolating disks for enclosed poles, local Laurent coefficient
certificates, and a finite winding-number computation.  These are exactly the
inputs that a classical diagram tends to suppress.

A uniform convergence theorem must include a primitive-recursive modulus.  A
pointwise statement without a modulus generally does not support interchange
of limit, integral, or derivative in \(\ipra\).  Dominated convergence is
replaced by finite truncation plus a tail or uniform-integrability modulus.

\subsection{Elementary and Diophantine number theory}

Bounded trial division, Euclid's algorithm, extended gcd, modular arithmetic,
finite factorization, valuations below a supplied bound, and Chinese remainder
construction are primitive-recursive finite algorithms.  Their proofs use
state invariants such as
\[
 r_i=A u_i+B v_i,
 \qquad
 \gcd(r_i,r_{i+1})=\gcd(A,B),
\]
encoded by integer equations and a bounded decreasing measure.

Diophantine geometry adds a search interface:
\begin{itemize}
\item a decidable point predicate;
\item a height and enumeration of bounded-height codes;
\item a height, descent, Runge, or covering certificate;
\item a finite verifier proving that the certificate applies to the input.
\end{itemize}
The generic descent theorem can then be fully embedded even when the production
of a global height certificate remains open.

\subsection{Analytic number theory}

Analytic-number-theory proofs should be divided into a finite arithmetic core
and an analytic input contract.

\begin{definition}[Analytic input certificate]
An analytic input certificate is a finite code containing explicit constants,
regions, exceptional cases, truncation lengths, and rational inequalities
sufficient for the later finite estimate.  Its checker validates every finite
algebraic and numerical condition.
\end{definition}

For a zero-free region, for example, the certificate records the region
parameters, a finite cover or local argument data, bounds on coefficients and
tails, and any exceptional zero with a separation estimate.  The downstream
prime-counting proof is a certified transducer from this data to an explicit
threshold.  The theorem is valuable even before a uniform certificate producer
is known: it states exactly which quantitative input is missing and prevents
an ineffective existence statement from entering unnoticed.

\subsection{Finite and effective probability}

A finite rational law is a finite list of rational values and natural weights
with positive total weight.  Expectation is an exact finite weighted sum.  A
product law is a mixed-radix enumeration; independence is an identity of
finite sums rather than an axiom.

Moments, characteristic functions at rational arguments, convolution, and
events with rational thresholds are finite calculations.  A Berry--Esseen or
Lindeberg proof is then a chain of polynomial/trigonometric approximation
certificates, finite product telescoping, and rational smoothing estimates.

An infinite or continuous law is admitted only through finite rational
approximants with explicit tail and quantization moduli.  For weak convergence,
a bounded-Lipschitz or characteristic-function metric is usually more robust
than CDF comparison.  Transfer to CDFs requires a continuity or
anti-concentration modulus near the threshold.  This cleanly separates the
finite probabilistic theorem from the presentation of the input law.

\subsection{Combinatorics and topology}

Finite combinatorics is a natural fit: graphs, words, partitions, tableaux,
simplicial complexes, and finite colorings are codes, while recurrences and
injections are bounded traces.  Generating-function proofs should be compiled
to coefficient identities before analytic notation is introduced.

Constructive topology uses subdivisions and parity certificates.  Sperner,
Brouwer-style approximate fixed points, polygonal Jordan separation, and
finite-dimensional degree arguments can be stated on rational grids with
explicit moduli.  Exact points arise from nested cells only when coherence is
proved.  Dimension and invariance results require careful negative statements,
but their local combinatorial certificates still fit finite-trace reflection.
